\documentclass[journal,onecolumn]{IEEEtran}

% =========================================================
% 宏包导入
% =========================================================
\usepackage[UTF8]{ctex}       % 核心：中文支持包
\usepackage{cite}             % 参考文献引用
\usepackage{amsmath,amssymb,amsfonts} % 数学公式支持
\usepackage{algorithmic}      % 算法伪代码
\usepackage{graphicx}         % 图片支持
\usepackage{textcomp}
\usepackage{xcolor}
\usepackage{bm}               % 粗体数学符号
\usepackage{algorithm}        % 算法环境
\usepackage{subfigure}        % 子图支持
\usepackage{url}

% =========================================================
% 自定义数学命令
% =========================================================
\newcommand{\norm}[1]{\left\lVert#1\right\rVert} % 范数符号 ||x||
\newcommand{\herm}{H} % 共轭转置符号 H
\newcommand{\trans}{T} % 转置符号 T
\newcommand{\trace}{\text{Tr}} % 迹

\begin{document}
	
	% =========================================================
	% 论文标题
	% =========================================================
	\title{基于多分支架构 PPO 的无人机辅助\\ RSMA 网络能效优化轨迹与资源联合分配}
	
	% 作者信息
	\author{ 
	}
	
	\maketitle
	
	
	
	\section{系统模型与问题建模}
	\label{sec:system_model}
	
	本文考虑一个下行空地（Air-to-Ground, A2G）通信系统，其中部署了一架旋翼无人机（UAV）以为 $K$ 个单天线静止地面用户（Ground Users, GU）提供服务，用户集合记为 $\mathcal{K} = \{1, 2, \dots, K\}$。为了充分利用三维空间的自由度，UAV 搭载了包含 $M = N_x \times N_y$ 根天线的均匀矩形阵列（Uniform Rectangular Array, URA）。我们将总飞行任务时间 $T$ 离散化为 $N$ 个等长的时隙，即 $T = N \delta_t$，其中 $\delta_t$ 为时隙长度。假设信道经历块衰落（Block Fading），即在单个时隙内信道保持恒定，而在不同时隙间发生变化。
	
	\subsection{无人机运动学与能耗模型}
	假设 UAV 在固定高度 $H$ 飞行。在第 $n$ 个时隙，UAV 的水平轨迹坐标记为 $\mathbf{q}[n] = [x[n], y[n]]^T \in \mathbb{R}^{2 \times 1}$。静止地面用户的位置已知且固定，记为 $\mathbf{w}_k = [x_k, y_k]^T, \forall k \in \mathcal{K}$。UAV 的飞行受限于其最大飞行速度 $V_{\max}$，其运动学更新方程可表示为：
	\begin{align}
		\mathbf{q}[n+1] &= \mathbf{q}[n] + \mathbf{v}[n] \delta_t, \quad \forall n \label{eq:kinematics} \\
		||\mathbf{v}[n]|| &\leq V_{\max} \label{eq:v_max}
	\end{align}
	其中，$\mathbf{v}[n]$ 为 UAV 在时隙 $n$ 的水平速度矢量。
	
	对于旋翼无人机，其总能耗由通信能耗和推进能耗组成。根据经典的旋翼无人机能耗分析模型，推进功率 $P_{fly}$ 作为飞行速度 $V = ||\mathbf{v}[n]||$ 的函数可表示为：
	\begin{equation}
		P_{fly}(V) = P_0 \left(1 + \frac{3 V^2}{U_{tip}^2} \right) + P_i \left( \sqrt{1 + \frac{V^4}{4 v_0^4}} - \frac{V^2}{2 v_0^2} \right)^{\frac{1}{2}} + \frac{1}{2} d_0 \rho s A V^3
		\label{eq:propulsion_power}
	\end{equation}
	其中，$P_0$ 和 $P_i$ 分别表示悬停状态下的桨叶型阻功率和诱导功率。其余参数（$U_{tip}, v_0, d_0, \rho, s, A$）均为与无人机空气动力学相关的常数。因此，UAV 在时隙 $n$ 的总能耗可表示为：
	\begin{equation}
		E_{tot}[n] = \left( P_{fly}(||\mathbf{v}[n]||) + P_{tx}[n] \right) \delta_t
		\label{eq:total_energy}
	\end{equation}
	其中，$P_{tx}[n]$ 为时隙 $n$ 的总发射功率。
	
	\paragraph{实现说明（简化能耗模型）}
	为降低训练开销并提升策略迭代稳定性，本文工程实现采用简化推进功率模型
	\begin{equation}
		P_{fly}^{\text{sim}}(V)=P_{\text{base}}+k_v V^2,
	\end{equation}
	并与通信功率 $P_{tx}[n]$ 共同构成时隙总功耗。该近似仅用于强化学习训练与对比实验，不改变本文的优化目标形式与约束建模。
	
	\subsection{基于 URA 的空地莱斯信道模型}
	在时隙 $n$，UAV 与第 $k$ 个用户之间的三维距离为 $d_k[n] = \sqrt{H^2 + ||\mathbf{q}[n] - \mathbf{w}_k||^2}$。由于 A2G 信道通常由视距（Line-of-Sight, LoS）链路主导，同时伴随非视距（NLoS）多径散射，本文采用莱斯（Rician）衰落信道模型。信道矢量 $\mathbf{h}_k[n] \in \mathbb{C}^{M \times 1}$ 建模为：
	\begin{equation}
		\mathbf{h}_k[n] = \sqrt{\beta_0 (d_k[n])^{-\alpha}} \left( \sqrt{\frac{K_R}{K_R+1}} \mathbf{a}(\phi_k[n], \theta_k[n]) + \sqrt{\frac{1}{K_R+1}} \mathbf{g}_k[n] \right)
		\label{eq:channel}
	\end{equation}
	其中，$\beta_0$ 为参考距离处的信道增益，$\alpha$ 为路径损耗指数，$K_R$ 为莱斯因子。矢量 $\mathbf{g}_k[n] \sim \mathcal{CN}(\mathbf{0}, \mathbf{I}_M)$ 表示瑞利衰落分量。
	
	假设 UAV 在飞行过程中保持水平姿态稳定，URA 的导向矢量 $\mathbf{a}(\phi_k[n], \theta_k[n])$ 完全由几何关系决定。离开方位角（Azimuth AoD）$\phi_k[n]$ 和离开俯仰角（Elevation AoD）$\theta_k[n]$ 可计算为：
	\begin{equation}
		\phi_k[n] = \arctan \left( \frac{y_k - y[n]}{x_k - x[n]} \right), \quad \theta_k[n] = \arcsin \left( \frac{H}{d_k[n]} \right)
		\label{eq:angles}
	\end{equation}
	此处arctan泛指考虑了四个象限的四象限反正切函数（Four-quadrant inverse tangent）。
	
	\subsection{RSMA 信号传输模型}
	本文采用单层 RSMA（1-layer RSMA）传输方案。在时隙 $n$，基站将 $K$ 个用户的消息拆分并重新编码为一个公共流 $s_c[n]$ 和 $K$ 个私有流 $s_k[n]$。发射信号矢量 $\mathbf{x}[n] \in \mathbb{C}^{M \times 1}$ 表示为：
	\begin{equation}
		\mathbf{x}[n] = \mathbf{p}_c[n] s_c[n] + \sum_{k=1}^K \mathbf{p}_k[n] s_k[n]
		\label{eq:transmit_signal}
	\end{equation}
	其中，$\mathbf{p}_c[n]$ 和 $\mathbf{p}_k[n]$ 分别为公共流和私有流的预编码矢量。总发射功率受限于最大功率约束，即 $P_{tx}[n] = ||\mathbf{p}_c[n]||^2 + \sum_{k=1}^K ||\mathbf{p}_k[n]||^2 \leq P_{\max}$。
	
	在接收端，假设发送端具备完美的瞬时信道状态信息（Perfect CSIT）且接收端能实现完美的串行干扰消除（Perfect SIC）。用户 $k$ 首先将所有私有流视为噪声来解码公共流 $s_c[n]$，其接收信号的信干噪比（SINR）为：
	\begin{equation}
		\gamma_{c,k}[n] = \frac{|\mathbf{h}_k^H[n] \mathbf{p}_c[n]|^2}{\sum_{j=1}^K |\mathbf{h}_k^H[n] \mathbf{p}_j[n]|^2 + \sigma^2}
		\label{eq:sinr_c}
	\end{equation}
	其中 $\sigma^2$ 为加性高斯白噪声的功率。为了确保所有用户均能成功解码公共流，公共流的可达速率必须受限于信道条件最差的用户，即 $R_c[n] = \min_{k \in \mathcal{K}} \log_2(1 + \gamma_{c,k}[n])$。
	
	在成功解码并利用 SIC 消除公共流后，用户 $k$ 继续解码自身的私有流 $s_k[n]$，其 SINR 为：
	\begin{equation}
		\gamma_k[n] = \frac{|\mathbf{h}_k^H[n] \mathbf{p}_k[n]|^2}{\sum_{j \neq k}^K |\mathbf{h}_k^H[n] \mathbf{p}_j[n]|^2 + \sigma^2}
		\label{eq:sinr_k}
	\end{equation}
	私有流的可达速率为 $R_k[n] = \log_2(1 + \gamma_k[n])$。设 $c_k[n]$ 为分配给用户 $k$ 的公共速率部分，其需满足 $\sum_{k=1}^K c_k[n] \leq R_c[n]$ 且 $c_k[n] \geq 0$。因此，用户 $k$ 在时隙 $n$ 的总可达速率为：
	\begin{equation}
		R_{k,tot}[n] = c_k[n] + R_k[n]
		\label{eq:total_rate}
	\end{equation}
	
	\subsection{优化问题构建}
	本文的优化目标是联合设计 UAV 的飞行轨迹 $\mathbf{V} \triangleq \{\mathbf{v}[n]\}$、RSMA 预编码矢量 $\mathbf{P} \triangleq \{\mathbf{p}_c[n], \mathbf{p}_k[n]\}$ 以及公共速率分配策略 $\mathbf{C} \triangleq \{c_k[n]\}$，以最大化整个飞行任务期间系统的比例公平能效（Proportional-Fair Energy Efficiency, PF-EE）。本文不再采用“公平硬约束优先”的建模思路，而是通过 PF-EE 分子中的对数效用对低速率用户赋予更高的边际收益权重，从而在统一目标中实现公平与能效的软权衡。该优化问题可建模为：
	\begin{align}
		\text{(P1)}: \max_{\mathbf{V}, \mathbf{P}, \mathbf{C}} \quad & \eta_{PF} = \frac{\sum_{n=1}^N \sum_{k=1}^K \log\left(1 + R_{k,tot}[n]\right)}{\sum_{n=1}^N \left( P_{fly}(||\mathbf{v}[n]||) + ||\mathbf{p}_c[n]||^2 + \sum_{k=1}^K ||\mathbf{p}_k[n]||^2 \right)} \label{eq:objective} \\
		\text{s.t.} \quad 
		& \text{(C1): } \mathbf{q}[n+1] = \mathbf{q}[n] + \mathbf{v}[n] \delta_t, \quad \forall n \tag{\ref{eq:objective}a} \\
		& \text{(C2): } ||\mathbf{v}[n]|| \leq V_{\max}, \quad \forall n \tag{\ref{eq:objective}b} \\
		& \text{(C3): } ||\mathbf{p}_c[n]||^2 + \sum_{k=1}^K ||\mathbf{p}_k[n]||^2 \leq P_{\max}, \quad \forall n \tag{\ref{eq:objective}c} \\
		& \text{(C4): } \sum_{k=1}^K c_k[n] \leq \min_{j \in \mathcal{K}} \log_2(1 + \gamma_{c,j}[n]), \quad \forall n \tag{\ref{eq:objective}d} \\
		& \text{(C5): } c_k[n] \geq 0, \quad \forall k, n \tag{\ref{eq:objective}e}
	\end{align}
	
	其中，目标函数分子为全局对数效用和，用于刻画比例公平准则；分母为任务总能耗。由于目标函数仍呈分式形式，且三维空间轨迹与空域预编码矢量之间存在强烈的非线性耦合，问题 (P1) 依然是一个高度非凸的非线性规划问题。传统的数学优化方法（如块坐标下降法或连续凸逼近）在求解此类问题时计算复杂度极高且难以避免近似误差。因此，在下一节中，我们将 (P1) 重新建模为马尔可夫决策过程（MDP），并提出一种新颖的多分支近端策略优化（MB-PPO）框架来高效求解该联合优化问题。
	
	
	
	\section{基于多分支深度强化学习的联合优化框架}
	\label{sec:algorithm}
	
	由于原优化问题 (P1) 是一个非凸、强耦合的分式优化问题，且比例公平目标与轨迹控制、空域预编码之间存在深度耦合，传统的交替优化算法难以在多维连续空间中实现低复杂度求解。因此，本节将原问题重构为马尔可夫决策过程（MDP），并提出一种新颖的多分支近端策略优化（Multi-Branch Proximal Policy Optimization, MB-PPO）算法。该框架通过定制化的动作空间映射，严格保证了物理约束的满足，并有效消除了多维异构动作更新时的梯度干扰。
	
	\subsection{马尔可夫决策过程 (MDP) 建模}
	MDP 由状态空间 $\mathcal{S}$、动作空间 $\mathcal{A}$ 和奖励函数 $\mathcal{R}$ 定义。在每个时隙 $n$，部署在 UAV 上的智能体观察当前状态 $s[n] \in \mathcal{S}$，采取动作 $a[n] \in \mathcal{A}$，并获得环境反馈的奖励 $r[n] \in \mathcal{R}$，同时状态转移至 $s[n+1]$。
	
	\subsubsection{状态空间设计 (State Space)}
	为了使智能体能够感知时变的空地信道、UAV 空间位置以及用户的公平性需求，我们设计了一种包含“速率历史感知”的异构状态表示。定义针对用户 $k$ 的局部特征向量为 $\mathbf{o}_k[n]$，其包含：
	a) 信道状态信息 (CSI)：将复数信道向量拆分为实部与虚部，即 $\Re(\mathbf{h}_k[n])$ 和 $\Im(\mathbf{h}_k[n])$。
	
	b) 空间位置信息：工程实现中状态输入采用\textbf{相对位置}而非绝对坐标，即 UAV-用户相对向量（例如 $\mathbf{w}_k-\mathbf{q}[n]$）及其归一化表示，以增强平移不变性。
	
	c) 公平性感知特征：为了引导网络向“饥饿”用户倾斜资源，引入上一时隙的可达速率 $R_{k,tot}[n-1]$。在比例公平目标下，对数效用会自动提升低速率用户的边际收益，从而强化公平导向。
	因此，时隙 $n$ 的全局状态表示为所有用户特征拼接而成的向量，即 $s[n] = \{ \mathbf{o}_1[n], \mathbf{o}_2[n], \dots, \mathbf{o}_K[n] \}$。
	
	\subsubsection{动作空间与物理约束映射 (Action Space and Constraint Mapping)}
	为了解决不同物理量纲和约束条件带来的挑战，我们采用全连续动作空间，并通过确定的代数映射将神经网络输出域 $[-1, 1]$ 严格投影至可行域内。定义时隙 $n$ 的动作 $a[n]$ 包含以下四个子集：
	\begin{itemize}
		\item \textbf{UAV 运动控制}：网络输出 $a_v[n], a_\theta[n] \in [-1, 1]$。UAV 的飞行速度和航向角通过线性映射获得：
		\begin{equation}
			v[n] = \frac{V_{\max}}{2}(a_v[n] + 1), \quad \theta[n] = \pi(a_\theta[n] + 1)
			\label{eq:action_traj}
		\end{equation}
		据此可计算二维速度向量 $\mathbf{v}[n] = [v[n]\cos\theta[n], v[n]\sin\theta[n]]^T$，严格满足约束 (C2)。
		
		\item \textbf{公共流方向分支}：网络输出公共流未归一化实数向量，重构为复向量后进行 $\ell_2$ 归一化，得到单位方向矢量 $\mathbf{\bar{p}}_c[n] \in \mathbb{C}^{M \times 1}$，满足 $||\mathbf{\bar{p}}_c[n]||=1$。
		
		\item \textbf{私有流方向分支}：网络输出 $K$ 个私有流未归一化实数向量，逐流重构并进行 $\ell_2$ 归一化，得到单位方向矢量 $\mathbf{\bar{p}}_k[n] \in \mathbb{C}^{M \times 1},\forall k$，满足 $||\mathbf{\bar{p}}_k[n]||=1$。
		
		\item \textbf{资源权重分支（$\rho,\alpha,\beta$）}：该分支联合输出总功率缩放、功率分配与公共速率分配。第一个变量 $a_{\rho}[n]\in[-1,1]$ 通过线性映射得到
		$\rho[n]=\frac{1}{2}(a_{\rho}[n]+1)$，对应总发射功率 $P_{tx}[n]=\rho[n]P_{\max}$；随后 $K+1$ 个变量经 Softmax 得到功率分配系数 $\alpha_c[n],\alpha_k[n]$，满足 $\alpha_c[n]+\sum_{k=1}^K\alpha_k[n]=1$；最后 $K$ 个变量经 Softmax 得到公共速率分配比例 $\beta_k[n]$，满足 $\sum_{k=1}^K\beta_k[n]=1$。因此预编码重构为：
		\begin{equation}
			\mathbf{p}_c[n] = \sqrt{\rho[n] \alpha_c[n] P_{\max}} \mathbf{\bar{p}}_c[n], \quad \mathbf{p}_k[n] = \sqrt{\rho[n] \alpha_k[n] P_{\max}} \mathbf{\bar{p}}_k[n]
			\label{eq:action_precoding}
		\end{equation}
		该联合映射在动作层面显式实现“方向-功率-速率”解耦，同时严格满足发射功率与分配比例约束。公共流分配速率为 $c_k[n]=\beta_k[n]R_c[n]$，满足约束 (C4) 与 (C5)。
	\end{itemize}
	
	\subsubsection{奖励函数设计 (Reward Function)}
	为与问题 (P1) 保持一致，本文将即时奖励定义为瞬时比例公平能效：
	\begin{equation}
		r[n] = \eta_{PF}[n] = \frac{\sum_{k=1}^K \log\left(1 + R_{k,tot}[n]\right)}{P_{fly}(||\mathbf{v}[n]||) + P_{tx}[n]}
		\label{eq:reward}
	\end{equation}
	其中，$\eta_{PF}[n]$ 为时隙 $n$ 的瞬时比例公平能效。由于对数效用对低速率用户具有更高的边际增益，智能体在优化系统能效的同时能够自然兼顾用户间公平性。需要说明的是，式 (\ref{eq:reward}) 属于对回合级分式目标 (P1) 的逐时代理优化信号，二者在一般情形下并非严格等价；本文通过确定性评估阶段的 PF-EE 与公平指标联合统计来验证该代理目标的有效性。
	
	\paragraph{实现说明（固定惩罚系数）}
	工程实现中采用固定权重的辅助惩罚项，以维持几何与物理可行性。实现奖励写为
	\begin{equation}
		r[n]=\eta_{PF}[n]-\omega_{core}\zeta_{core}[n]-\zeta_{wall}[n]-\zeta_{pow}[n].
		\label{eq:reward_impl_20260415}
	\end{equation}
	其中 $\omega_{core}>0$ 为常数权重。由于本文主目标已通过对数效用显式刻画公平性，固定惩罚系数即可稳定引导可行域探索，无需额外退火调度。
	
	
	
	\subsubsection{比例公平目标驱动机制}
	本文采用无固定阈值的目标驱动训练方式，不再引入额外的速率违规对偶变量，也不设置 Jain 指数或最弱用户速率的硬下界约束。Actor 与 Critic 直接围绕式 (\ref{eq:reward}) 所定义的比例公平奖励进行联合优化，以学习公平与能效之间的帕累托权衡关系。工程惩罚项 $\zeta_{core}[n]$、$\zeta_{wall}[n]$ 与 $\zeta_{pow}[n]$ 仅用于保持几何与物理可行性，不承担阈值约束更新角色。该设计降低了训练复杂度，并缓解了由附加调度机制引发的超参数敏感性。
	
	
	
	
	
	\subsection{多分支网络架构 (Multi-Branch Architecture)}
	在处理多维耦合动作（如缓慢变化的无人机轨迹与快速变化的空域预编码）时，传统的单头 Actor 网络极易面临严重的梯度干扰（Gradient Interference）问题，导致策略难以收敛。为此，我们设计了具有对齐结构的多分支 Actor-Critic 架构。
	
	如图 X 所示（\textit{注：此处在论文中需配合网络架构图}），Actor 网络首先通过一个共享的特征提取模块（如多层感知机）处理拼接后的全局用户状态信息 $s[n]$。提取出的高维联合特征随后分化为 4 个独立的分支（Heads）：轨迹控制分支、公共流方向分支、私有流方向分支和资源权重分支（联合输出 $\rho,\alpha,\beta$）。每个分支独立输出对应连续动作分布的高斯均值（Mean）和标准差的对数（Log-Std）。这种“底层特征共享，顶层动作解耦”的设计，有效隔离了异构动作空间的梯度回传干扰。Critic 网络则采用单头结构，输出对全局状态的价值评估 $V(s[n])$。
	
	\subsection{MB-PPO 训练算法}
	MB-PPO 基于 Actor-Critic 框架进行在线训练。定义 Actor 的策略为 $\pi_\theta(a|s)$，Critic 的价值函数为 $V_\phi(s)$。在每一轮训练中，智能体收集一批经验轨迹数据，并利用广义优势估计（GAE, Generalized Advantage Estimation）计算优势函数 $A^{\pi}[n]$。为了保证策略更新的单调递增并防止更新步长过大破坏稳定性，MB-PPO 采用裁剪的代理目标函数（Clipped Surrogate Objective）更新策略网络参数 $\theta$：
	\begin{equation}
		L^{CLIP}(\theta) = \mathbb{E} \left[ \min(r_t(\theta) A^{\pi}[n], \text{clip}(r_t(\theta), 1-\epsilon, 1+\epsilon) A^{\pi}[n]) \right]
		\label{eq:ppo_clip}
	\end{equation}
	其中 $r_t(\theta) = \frac{\pi_\theta(a[n]|s[n])}{\pi_{\theta_{old}}(a[n]|s[n])}$ 为新旧策略的比率，$\epsilon$ 为裁剪超参数。Critic 网络参数 $\phi$ 则通过最小化价值预测均方误差进行更新。
	
	
	\section{仿真结果与讨论}
	\subsection{新目标函数下的评价指标定义}
	为与比例公平能效目标保持一致，本文在仿真中采用以下评价指标。为便于与式 (\ref{eq:objective}) 保持一致，下文默认使用“功率和”作为等价能耗指标。
	
	1) \textbf{比例公平能效（PF-EE）}：作为主指标，定义为
	\begin{equation}
		\eta_{PF}^{ep}=\frac{\sum_{n=1}^{N}\sum_{k=1}^{K}\log\left(1+R_{k,tot}[n]\right)}{\sum_{n=1}^{N}\left(P_{fly}(||\mathbf{v}[n]||)+P_{tx}[n]\right)}.
		\label{eq:metric_pf_ee}
	\end{equation}
	
	2) \textbf{传统能效（EE）}：用于与既有方法横向对照，定义为
	\begin{equation}
		\eta_{EE}^{ep}=\frac{\sum_{n=1}^{N}\sum_{k=1}^{K}R_{k,tot}[n]}{\sum_{n=1}^{N}\left(P_{fly}(||\mathbf{v}[n]||)+P_{tx}[n]\right)}.
		\label{eq:metric_ee}
	\end{equation}
	
	3) \textbf{Jain 公平指数}：用于刻画多用户速率均衡性。先定义每时隙公平指数
	\begin{equation}
		J[n]=\frac{\left(\sum_{k=1}^{K}R_{k,tot}[n]\right)^2}{K\sum_{k=1}^{K}R_{k,tot}^2[n]},
		\label{eq:metric_jain_slot}
	\end{equation}
	再取回合平均
	\begin{equation}
		\bar{J}^{ep}=\frac{1}{N}\sum_{n=1}^{N}J[n].
		\label{eq:metric_jain_ep}
	\end{equation}
	
	4) \textbf{最弱用户平均速率}：用于评估边缘用户服务质量，定义为
	\begin{equation}
		\bar{R}_{\min}^{ep}=\frac{1}{N}\sum_{n=1}^{N}\min_{k\in\mathcal{K}}R_{k,tot}[n].
		\label{eq:metric_rmin}
	\end{equation}
	
	5) \textbf{总能耗}：用于评估能量开销本身，定义为
	\begin{equation}
		E_{tot}^{ep}=\sum_{n=1}^{N}\left(P_{fly}(||\mathbf{v}[n]||)+P_{tx}[n]\right)\delta_t.
		\label{eq:metric_etot}
	\end{equation}

	6) \textbf{公平收益-能效代价比}：用于直接验证“牺牲少部分能效换取较大公平提升”。设参考策略（如 Sum-EE）为 ref、待比较策略为 alg，定义
	\begin{equation}
		L_{EE}=\frac{\eta_{EE,\mathrm{ref}}^{ep}-\eta_{EE,\mathrm{alg}}^{ep}}{\eta_{EE,\mathrm{ref}}^{ep}},\quad
		G_J=\frac{\bar{J}_{\mathrm{alg}}^{ep}-\bar{J}_{\mathrm{ref}}^{ep}}{\bar{J}_{\mathrm{ref}}^{ep}},\quad
		G_{R}=\frac{\bar{R}_{\min,\mathrm{alg}}^{ep}-\bar{R}_{\min,\mathrm{ref}}^{ep}}{\bar{R}_{\min,\mathrm{ref}}^{ep}},
		\label{eq:metric_tradeoff_gain_loss}
	\end{equation}
	并进一步定义
	\begin{equation}
		\Xi_J=\frac{G_J}{L_{EE}+\epsilon},\quad \Xi_R=\frac{G_R}{L_{EE}+\epsilon},
		\label{eq:metric_tradeoff_ratio}
	\end{equation}
	其中 $\epsilon>0$ 为防止分母为零的微小常数。当 $L_{EE}$ 保持在较小范围而 $\Xi_J$ 或 $\Xi_R$ 显著大于 1 时，可判定系统达成“少量 EE 损失对应较大公平增益”的权衡目标。
	
	上述指标共同反映“公平-能效-可行性”三方面性能：$\eta_{PF}^{ep}$ 作为主优化目标，$\bar{J}^{ep}$ 与 $\bar{R}_{\min}^{ep}$ 衡量公平收益，$\eta_{EE}^{ep}$ 与 $E_{tot}^{ep}$ 衡量能效与能耗代价，而 $\Xi_J$、$\Xi_R$ 用于量化公平提升相对于能效牺牲的性价比。

	\subsection{PF-EE 权衡命题的仿真验证准则}
	为验证本文命题“通过牺牲少部分能效较大幅度提升公平”，仿真采用统一场景与统一训练预算，对不同目标偏好的策略（如 Sum-EE、PF-EE、Max-Min）进行\textbf{确定性评估}并在尾部稳定区间统计均值。记参考策略为 ref（通常取 Sum-EE），若满足
	\begin{equation}
		L_{EE} \leq \tau_{EE},\quad G_J \geq \tau_J \ \text{或}\  G_R \geq \tau_R,
		\label{eq:tradeoff_criterion}
	\end{equation}
	则认为该策略在工程意义上实现了“以小代价换大公平收益”的目标。参数 $\tau_{EE},\tau_J,\tau_R$ 可按任务需求设定（例如可取小代价阈值与中高公平增益阈值），并在实验设置中统一报告。
	
	
\end{document}